% Research fragment for the main report. All repository references use the
% pinned revision c39c388d4f87220ab9e758ae457520bbc0d74341.
\section{The boundary between certificates and global witnesses}
\label{sec:certificate-boundary}

The audited arithmetic packets \cite{faltings,siegel,mordell,roth,prime} make an essential distinction that a general
embedding framework must preserve. A finite computation can check a
candidate point, a rational identity, or a finite formal derivation.
It cannot automatically certify that all points of an infinite search
space have already been found. The latter is a mathematical assertion
about every possible point. There is also a distinction between a proof
that some bound exists and a certificate which contains a particular
numerical bound and certifies its correctness.

\subsection{An exact impossibility theorem for complete bound certificates}

Fix an effective enumeration of deterministic programs. Let $H(e,t)$ mean
that program $e$, run on input $e$, halts for the first time at positive
step $t$. Bounded simulation and a bounded check of the preceding steps
make $H$ primitive recursive. Determinism implies
\[
 H(e,s)\wedge H(e,t)\ \longrightarrow\ s=t.
\]
Thus every fiber has the uniform cardinality bound one. Set
\[
 \operatorname{Bound}(e,B)
   \quad\Longleftrightarrow\quad
 \forall t\bigl(H(e,t)\longrightarrow t\le B\bigr).
\]

\begin{theorem}[No complete effective system of numerical bound certificates]
\label{thm:no-complete-bound-certificates}
There is no total computable predicate $C(e,B,w)$ satisfying both of the
following properties in the standard natural numbers:
\begin{enumerate}
\item \emph{Soundness:} $C(e,B,w)$ implies
      $\operatorname{Bound}(e,B)$.
\item \emph{Completeness:} for every $e$ there are $B,w$ with $C(e,B,w)$.
\end{enumerate}
The same impossibility holds for a computably enumerable acceptance
relation. In particular, it holds for primitive recursive checkers.
\end{theorem}

\begin{proof}
Suppose first that $C$ is decidable. Given $e$, enumerate all pairs
$(B,w)$ in a fixed effective order and test $C(e,B,w)$. Completeness
guarantees that this search terminates. Its running time need not have
a primitive recursive bound; ordinary recursive termination is already
enough for the argument. Let $B_e$ be the bound in the first accepted
certificate.

Simulate program $e$ on input $e$ for $B_e$ steps. If it halts, answer
yes. If it does not halt within those steps, answer no. Soundness makes
the second answer correct: any later first halting time $t$ would satisfy
$t>B_e$, contrary to $\operatorname{Bound}(e,B_e)$. Hence this procedure
would be a total decision algorithm $D$ for diagonal halting.

For completeness, the diagonal contradiction is immediate. Construct a
program $P$ which, on input $n$, runs $D(n)$ and loops forever when the
answer is yes, but halts when the answer is no. Let $p$ be a code of $P$.
If $D(p)$ answers yes, then $P(p)$ loops; if it answers no, then $P(p)$
halts. Either answer contradicts the asserted correctness of $D$.

If acceptance is only computably enumerable, dovetail the acceptance
computations for all pairs $(B,w)$. Completeness again supplies a pair
whose acceptance will eventually be observed, and the remainder of the
argument is unchanged.
\end{proof}

The theorem concerns certificates which explicitly supply a numerical
bound. It does not say that a classical theory cannot prove
$\forall e\exists B\operatorname{Bound}(e,B)$: classical case analysis
on whether $e(e)$ halts proves that assertion. A finite classical proof
of the existential sentence does not by itself provide a numerical
witness. Nor does the theorem say that any particular family of
algebraic curves has this obstruction. The halting fibers are a
counterexample to a general certificate-completeness principle, not
an encoding of curves.

This yields a practical requirement for the framework. Every use of
certificates should identify separately the finite acceptance relation,
its proved soundness, and the existence or construction of accepted
certificates on the intended input class. A sound checker may be a useful
artifact even when completeness is unavailable. It cannot serve as a
substitute for the missing completeness theorem.

\subsection{The precise intuitionistic witnessing theorem}

The strong witnessing assertion used in the Siegel packet is supported
by a primary proof-theoretic reference. Salehi defines $\mathsf{iPRA}$
using all primitive recursive defining equations and atomic induction,
and states the following as Fact~3.5 on page~57:
\[
 \begin{gathered}
 T\vdash\forall x\exists y\,A(x,y)\\
 \Longrightarrow\quad
 \text{some primitive recursive symbol }f\text{ satisfies}\\
 T\vdash\forall x\,A(x,f(x)).
 \end{gathered}
\]
for $T=\mathsf{iPRA}$ or intuitionistic $\Sigma_1$ induction in the
primitive recursive language. The matrix $A$ is not restricted to a
decidable or quantifier-free formula. See \cite{salehi}.
This is an external metatheorem, not an internal reflection axiom and
not a result formalized by the present files. Atomic and quantifier-free
induction are interderivable here by the Boolean representations of
quantifier-free formulas.

Ferreira and Marques give a related formulation for intuitionistic
$I\Sigma_1$, with correctness in its primitive-recursive definitional
extension; their author manuscript states the theorem on printed
page~14:
\cite{ferreira}.
The Faltings packet cites that more conservative formulation and does
not claim an extracted correctness proof in $\mathsf{iPRA}$ itself.
Salehi's exact statement supports the stronger conclusion when its
specified $\mathsf{iPRA}$ hypothesis is met.

Here is the premise-management step in full. Suppose $G(x)$ is decidable
in $\mathsf{iPRA}$ and
\[
 \mathsf{iPRA}\vdash
 \forall x\bigl(G(x)\longrightarrow\exists y\,A(x,y)\bigr).
\]
Fix $x$. If $G(x)$, the displayed theorem supplies $y$ and $A(x,y)$,
and therefore $G(x)\to A(x,y)$. If $\neg G(x)$, choose $y=0$;
the implication $G(x)\to A(x,0)$ follows by contradiction. Decidability
of $G$ proves these cases exhaustive. Generalization gives
\[
 \mathsf{iPRA}\vdash
 \forall x\exists y\bigl(G(x)\longrightarrow A(x,y)\bigr).
\]
The witnessing theorem now supplies one primitive recursive $f$ and
an $\mathsf{iPRA}$ proof of
\[
 \forall x\bigl(G(x)\longrightarrow A(x,f(x))\bigr).
\]
This does not use unrestricted independence of premise. Its essential
hypothesis is the proved decidability of $G$.

Apply this observation with
$A(c,B)=\forall p(P(c,p)\to h(p)\le B)$. A uniform positive
$\mathsf{iPRA}$ theorem on decidably admissible curve codes gives a
primitive recursive height bound, even though $A$ contains an unbounded
universal quantifier. If $c$ includes a certificate, the extracted
function is allowed to depend on that certificate. A further producer
is necessary before the input can be reduced to raw curve coefficients.
The implication gives a necessary computational consequence; absence
of a known bound is not a nonprovability theorem.

\subsection{What the audited arithmetic packets still require}

The following statements are relative to the pinned repository revision,
not a claim that all mathematics in the packets has been verified.

\paragraph{Faltings and Siegel.}
The positive list and height versions require a global height bound.
The packets prove bounded enumeration after such a bound is supplied
and distinguish a cardinality bound from a height bound. The Siegel
packet already proves the finite covering last step of a quantitative
Subspace Theorem route: at most $M$ subspaces, each meeting the curve
in at most $D$ relevant points, give a cardinality bound $MD$. Its
missing ingredients include a finite-tuple approximation theorem,
the requisite curve-function presentations, lifting and valuation
estimates, and their arithmetic proofs. The present general framework
does not supply those ingredients.

\paragraph{Mordell--Weil.}
The supplied-cover result is substantially more explicit than an
abstract invocation of descent. For an integral short Weierstrass model
$y^2=x^3+Ax+B$, it gives a primitive recursive list of all halves of
any supplied rational point by a rational-root scan of a quartic. From
a supplied complete list of representatives modulo doubling it derives
an explicit height contraction and generator list. Pointwise
decidability of divisibility by two does not construct that complete
quotient-covering list. General production of that list remains an
unfilled obligation.

\paragraph{Roth.}
The quadratic case and the Liouville range are treated by explicit
polynomial separation. General cutoff effectivity is left open. The
packet already includes the cardinality-to-empty-window construction
using $N+1$ successive intervals, so that observation should not be
presented as a new contribution.

\paragraph{Chebotarev.}
The packet explicitly separates effective analytic estimates from their
translation to $\mathsf{iPRA}$. Its conditional density-modulus lemma
is supplied with integer thresholds. The still missing inputs include
number-field and prime-ideal algorithms with arithmetic correctness,
Dedekind/Hecke explicit formulas, quantitative zero-free estimates,
and certified primitive recursive constants. A possible exceptional
zero alone is not asserted to prevent effective Chebotarev.

\paragraph{Siegel--Walfisz and Bombieri--Vinogradov.}
The report isolates uniform constants and thresholds before the
unbounded size variable. Its elementary reduction from effective
Bombieri--Vinogradov witnesses to logarithmic-power Siegel--Walfisz
estimates shows why averaging alone does not erase the issue. It also
retains the aggregate rational approximation error $QX2^{-s}$ when
summing over $Q$ moduli. Neither finite evaluation of prime counts nor
a finite Gram-matrix certificate supplies the missing uniform analytic
estimates.

\subsection{Lean targets for the certificate boundary}

The computability obstruction is a metatheoretic theorem, naturally
formalized in Lean's ordinary metatheory rather than as a new
$\mathsf{iPRA}$ axiom. A clean dependency chain is:
\begin{enumerate}
\item Fix a concrete machine model and a primitive recursive bounded
simulation, with the first-halt predicate and its uniqueness proof.
\item Define $\operatorname{Bound}$ and a certificate structure with
an acceptance predicate and semantic soundness theorem.
\item Given totality of accepted pairs, define the partial-recursive
least accepted pair; prove its totality from the hypothesis. No claim
that this minimization is primitive recursive is needed.
\item Compose it with bounded simulation and prove that it decides
diagonal halting, using soundness exactly when the simulation reports
no halt by the supplied bound.
\item Apply the chosen machine model's diagonal undecidability
theorem, or formalize the short self-application contradiction above.
\end{enumerate}
For iPRA integration, only the finite simulator, first-halt uniqueness,
and finite acceptance checks are object-language kernels. The
impossibility theorem and the external witnessing theorem belong in
the metatheory and must not be added as arithmetic axioms.

